% !PRINTED 276
\weberpage{276}
It follows first of all from \S 69, (9):
\[
c_1 \equiv c+\partial \pmod{a}, \quad c_1' \equiv c'+\partial' \pmod{a'}, \quad c_1'' \equiv c''+\partial'' \pmod{a''},
\]
from which, by (4),
\[
\begin{aligned}
\partial''c_1' &\equiv \partial''c' + \partial \equiv c_1 \pmod{a'}\\
\partial'c_1'' &\equiv \partial'c'' + \partial \equiv c_1 \pmod{a''}
\end{aligned}
\]
in agreement with the congruences (4).

For the composition (6) it follows, by \S 69, (12), that
$\partial_2$, $\partial_2'$, $\partial_2''$ are the greatest common divisors of $a,c$;
$a',c'$; $a'',c''$; but since $\partial''$ is relatively prime to $a'$, and
$\partial''c' \equiv c \pmod{a'}$, $\partial_2'$ is also the greatest common
divisor of $a'$ and $c$, and for the same reasons $\partial_2''$ the greatest
common divisor of $a''$ and $c$. From this one concludes, since
$a'$, $a''$ are relatively prime:
\[
\partial_2 = \partial_2'\partial_2'', \qquad a_2 = a_2'a_2''.
\]
Further, by \S 69, (11), (12), (13):
\[
\begin{aligned}
\partial_2 &= \alpha c - \gamma a, & c_2 &= -\partial\alpha,\\
\partial_2' &= \alpha'c' - \gamma'a', & c_2' &= -\partial'\alpha',
\end{aligned}
\]
hence, by (4):
\[
\partial_2c_2' - c_2\partial_2' \equiv \alpha\alpha'(\partial c' - c\partial') \equiv 0 \pmod{n'},
\]
and consequently also
\[
\partial_2''c_2' - c_2 \equiv 0 \pmod{a_2'},
\]
in agreement with (4), and in the same way it follows that:
\[
\partial_2'c_2'' - c_2 \equiv 0 \pmod{a_2''},
\]
whereby the proof is furnished that the belonging together of the
transformations expressed by (4)
\[
\begin{pmatrix} a, & 0\\ c, & \partial \end{pmatrix}, \quad
\begin{pmatrix} a', & 0\\ c', & \partial' \end{pmatrix}, \quad
\begin{pmatrix} a'', & 0\\ c'', & \partial'' \end{pmatrix}
\]
is not disturbed by the application of any \gesperrt{linear} transformation to $\omega$.
Since we here assume $n$ to be odd, we may always take $c$, $c'$, $c''$
to be divisible by 16; and if $n$, and consequently also $n'$, $n''$, are
indivisible by 3, then $c$, $c'$, $c''$ may also be assumed divisible by 3.
But if $n$ is divisible by 3, then of the two factors $n'$, $n''$ one, say
$n''$, will be divisible by 3, and the other, $n'$, not. It may then still
be assumed that $c'$ is divisible by 3, but not $c$ and $c''$.
In this case the dependence of $c''$ upon $c$ shall be determined
more closely by the congruence
\weq{7}{
\partial'c'' \equiv c \pmod{3a''}.
}
